\scanpage{43}
% L1–3
groupes sur $\bar X$, également ind-fini en vertu de (IX 1.6. (ii)). \handwritten{\'}{E}crivons maintenant l'homomorphisme naturel de suites exactes d'ensembles pointés, où l\deleted{\unsure{es}}\handwritten{a} première flèche horizontale de chaque ligne est injective:\handwritten{\tikz[baseline=-.3ex]{\draw (0,.04) circle (.038) (0,.078)--(.03,.29);}}
\margin{\handwritten{\tikz[baseline=-.3ex]{\draw (0,.04) circle (.038) (0,.078)--(.03,.29);} réf}}

% L4–5
\diagram{\begin{tikzcd}[ampersand replacement=\&,column sep=small,row sep=small]
\deleted{\unsure{0}}\&H^1(X,\handwritten{\pi_*}(\bar F))\handwritten{\arrow[r]\arrow[d]}\&H^1(\bar X,\bar F)\handwritten{\arrow[r]\arrow[d]}\&H^0(X,R^1\handwritten{\pi_*}(\bar F))\handwritten{\arrow[d]}\\
\deleted{\unsure{0}}\&H^1(Y,\handwritten{\pi}_{Y\handwritten{*}}(\bar F_{\bar Y}))\handwritten{\arrow[r]}\&H^1(\bar Y,\bar F_{\handwritten{\bar Y}})\handwritten{\arrow[r]}\&H^0(Y,R^1\handwritten{\pi}_{Y\handwritten{*}}(\bar F_{\bar Y})),
\end{tikzcd}}

% L6–13
cf (\handwritten{\unsure{3?2.|3.2.}}). Compte tenu de l'hypothèse sur les homomorphismes de changement de base faites dans (i) et (ii), et de la co\deleted{\unsure{nditi}}\handwritten{nclusi}on de (i), la pre\handwritten{mi}ère flèche verticale s'identifie à la flèche de restriction $H^1(X,F)\handwritten{\to}H^1(\deleted{\unsure{X}}Y,F_Y)$ dont nous voulons établir la surjectivité, et la dernière flèche verticale est \handwritten{in}jective\handwritten{;} d'autre part, d'après l'hypothèse (ii), la deuxième flèche verticale est bijective. On en conclut aussitôt la \deleted{\unsure{bijectiv}} surjectivité de la première flèche verticale, ce qui achève de prouver (ii). \deleted{\unsure{En}}

% L14–17
(iii) \deleted{(\unsure{Esquisse|quisse} abrégée). On adopte pour abréger les notations des catégories dérivées (sinon, il faut passer à des résolutions injectives..). On considère l'homomorphisme canonique $F\handwritten{\to}R\handwritten{\pi_*}(\bar F)$, où $\bar F=\handwritten{\pi^*}(F)$, et son mapping-cylinder $G$, de sorte qu'on a un triangle exact}

% L18, triangle
\diagram{\begin{tikzcd}[ampersand replacement=\&,column sep=tiny,row sep=small]
\&\deleted{G}\handwritten{\arrow[dl]}\&\\
\deleted{F}\handwritten{\arrow[rr]}\&\&\deleted{R\handwritten{\pi_*}(\bar F)}\handwritten{\arrow[ul]}{\qquad\deleted{\text{.}}}
\end{tikzcd}}

% L19–23
\deleted{Pour \unsure{pouver|ppouver} \deleted{\unsure{le théorême}} l'assertion pour $F$, on est ramené à la prouver pour $R\handwritten{\pi_*}(\bar F)$ et $G$ séparément, grâce à la suite exacte de cohomologie et au lemme des cinq\deleted{\unsure{s}}. Pour $G$, on note que $G$ est \deleted{à cohomologie} \deleted{\unsure{faisceaux de}} nul\deleted{\unsure{le}} sur $U$, et pour un tel complexe de faiceaux, la suite spectrale de cohomologie nous ramène au cas d'un faisceau concentré sur $Y$, qui est trivial.}

% L24–25
\deleted{Le cas de $R\handwritten{\pi_*}(\bar F)$ se traite en notant que \deleted{\unsure{$\underline H R(\bar F)\quad\underline H$}} l'homomorphisme $\underline H^\bullet(X,R\handwritten{\pi_*}(\bar F))\handwritten{\to}H^\bullet(Y,h^{\handwritten{*}}(R\handwritten{\pi_*}(\bar F))$ s'identifie, grâce à l'hypothèse sur l'homo}

% L26–28, bottom in pen
\inserted{\handwritten{Adapter la démonstration de \uline{XII} 8.3. (iii); on se ramène si on veut au cas $\pi$ \uline{surjectif}, en remplaçant $\bar X$ par $\bar X\amalg Y$ !}}

% archive sheet number: 42
